\chapter[Marco Conceptual]{Marco Conceptual}
\label{ch:marco-conceptual}

\insertminitoc
\parindent0pt

\clearpage

\section{Glosario conceptual del estudio}

\begin{itemize}
    \item[] \textbf{Algoritmo de Recomendación:} Sistema de filtrado de información diseñado para predecir las preferencias de un usuario y sugerir contenido relevante, frecuentemente utilizado en redes sociales para maximizar el tiempo de retención (\textit{engagement}). \\
    
    \item[] \textbf{Arquitectura Limpia (\textit{Clean Architecture}):} Patrón de diseño de software que separa los elementos de un sistema en capas con reglas de dependencia estrictas, aislando la lógica de negocio de la interfaz de usuario y las bases de datos para facilitar su mantenimiento y escalabilidad. \\
    
    \item[] \textbf{Bitácora de Evidencias:} Registro digital estructurado y seguro donde se almacenan datos, eventos o capturas generadas por el sistema de monitoreo, con el fin de auditar incidentes de riesgo y facilitar la revisión parental. \\
    
    \item[] \textbf{Ciberacoso (\textit{Cyberbullying}):} Uso de medios de comunicación digital para acosar, amenazar o humillar a un individuo de forma deliberada y repetida, aprovechando a menudo el anonimato de la red. \\
    
    \item[] \textbf{Cifrado de Extremo a Extremo (\textit{End-to-End Encryption, E2EE}):} Método de comunicación segura que evita que terceros accedan a los datos mientras se transfieren de un sistema o dispositivo a otro, garantizando que solo el emisor y el receptor posean las claves criptográficas para leerlos. \\
    
    \item[] \textbf{Computación en el Borde (\textit{Edge Computing}):} Paradigma informático que desplaza el procesamiento de datos, la inferencia de Inteligencia Artificial y el almacenamiento lo más cerca posible de la fuente generadora (el dispositivo móvil), reduciendo la latencia y la dependencia de servidores en la nube. \\
    
    \item[] \textbf{Desinhibición Digital:} Alteración psicológica por la cual las personas se comportan en el ciberespacio con menor contención y empatía que en interacciones físicas, impulsada por la falta de señales no verbales y la percepción de anonimato. \\
    
    \item[] \textbf{Fatiga de Alerta (\textit{Alert Fatigue}):} Fenómeno de desensibilización que experimenta un usuario (como un padre o tutor) al recibir un volumen excesivo de notificaciones o falsos positivos de un sistema de seguridad, llevándolo a ignorar advertencias reales. \\
    
    \item[] \textbf{FoMO (\textit{Fear Of Missing Out}):} Ansiedad social caracterizada por el temor constante a que otros estén viviendo experiencias gratificantes de las cuales uno está ausente, fuertemente exacerbada por el uso continuo de redes sociales. \\
    
    \item[] \textbf{Grooming (Ciberengaño pederasta):} Proceso mediante el cual un adulto establece un vínculo de confianza emocional con un menor a través de internet con el objetivo último de someterlo a abuso o explotación sexual. \\
    
    \item[] \textbf{Humano en el Ciclo (\textit{Human-in-the-loop, HITL}):} Modelo de diseño de Inteligencia Artificial donde el sistema automatizado requiere la interacción, supervisión o validación humana antes de tomar una decisión o ejecutar una acción crítica. \\
    
    \item[] \textbf{Inteligencia Artificial Explicable (\textit{Explainable AI, XAI}):} Conjunto de procesos y métodos que permiten a los usuarios humanos comprender y confiar en los resultados o decisiones generadas por algoritmos de aprendizaje automático, evitando el fenómeno de "caja negra". \\
    
    \item[] \textbf{Mediación Habilitadora (\textit{Enabling Mediation}):} Estilo de crianza digital en el que los padres fomentan el uso seguro y responsable de la tecnología a través de la comunicación, la guía y el co-uso, en contraposición a la prohibición absoluta. \\
    
    \item[] \textbf{Privacidad por Diseño (\textit{Privacy by Design}):} Marco de trabajo ético y técnico que integra la protección de datos personales en todas las etapas de la ingeniería y diseño de un sistema informático, desde su concepción hasta su despliegue. \\
    
    \item[] \textbf{Procesamiento de Lenguaje Natural (\gls{nlp}):} Rama de la Inteligencia Artificial que dota a las computadoras de la capacidad de comprender, interpretar y manipular el lenguaje humano, utilizada para analizar el tono o la intención en los mensajes de texto. \\
    
    \item[] \textbf{Resiliencia Digital:} Capacidad de un individuo, especialmente un menor, para navegar por internet de forma segura, identificar amenazas, sobreponerse a experiencias negativas en línea y aprender de ellas. \\
    
    \item[] \textbf{Sexting:} Acrónimo de *sex* y *texting*; práctica consistente en el envío, recepción o reenvío de mensajes, fotografías o videos de contenido sexualmente explícito a través de dispositivos móviles. \\
    
    \item[] \textbf{Tecnoferencia:} Interrupción de las interacciones sociales y los vínculos afectivos cotidianos (por ejemplo, entre padres e hijos) debido a la atención intrusiva que demandan los dispositivos tecnológicos. \\
    
    \item[] \textbf{Uso Problemático de Internet (UPI):} Patrón de comportamiento caracterizado por una incapacidad para controlar el tiempo de uso de tecnologías digitales, provocando disfuncionalidad en la vida diaria, síntomas de abstinencia y aislamiento social. \\
    
    \item[] \textbf{Visión Artificial (\textit{Computer Vision}):} Campo de la Inteligencia Artificial que entrena a las computadoras para extraer información significativa y clasificar contenido a partir de imágenes digitales y videos en tiempo real, como el análisis de capturas de pantalla.

    \item[] \textbf{Aplicaciones Bóveda (\textit{Vault Apps}):} Aplicaciones móviles diseñadas con interfaces engañosas (como calculadoras o calendarios) cuyo propósito real es ocultar galerías, mensajes y datos cifrados fuera del radar de padres o tutores. \\

    \item[] \textbf{Aprendizaje Federado (\textit{Federated Learning}):} Enfoque de aprendizaje automático descentralizado donde un modelo de \gls{ia} se entrena localmente en múltiples dispositivos sin enviar los datos privados del usuario a la nube, intercambiando únicamente actualizaciones matemáticas. \\

    \item[] \textbf{Arquitectura Basada en Eventos:} Patrón de diseño de software donde el flujo del programa está determinado por eventos (como mensajes entrantes o cambios de pantalla), crucial para sistemas de monitoreo reactivos en tiempo real. \\

    \item[] \textbf{Cuantización de Modelos:} Técnica de optimización de Inteligencia Artificial que reduce la precisión de los parámetros matemáticos de una red neuronal (ej. de 32 a 8 bits) para permitir su ejecución fluida en procesadores de dispositivos móviles con batería limitada. \\

    \item[] \textbf{Desarrollo Positivo Juvenil (PYD):} Enfoque de la psicología evolutiva centrado en potenciar las fortalezas, competencias y la agencia del menor para construir resiliencia, en lugar de enfocarse únicamente en la represión de conductas de riesgo. \\

    \item[] \textbf{Diseño de Fricción (\textit{Friction Design}):} Estrategia de Interacción Humano-Computadora (HCI) que introduce intencionalmente pequeñas barreras o pasos adicionales en una interfaz para interrumpir el uso compulsivo o automatizado por parte del usuario. \\

    \item[] \textbf{Huella Digital (\textit{Digital Footprint}):} Rastro persistente de datos, metadatos y registros de actividad que un menor deja de forma activa o pasiva al navegar por internet y utilizar plataformas digitales, el cual puede ser perfilado comercialmente. \\

    \item[] \textbf{OCSEA (\textit{Online Child Sexual Exploitation and Abuse}):} Acrónimo internacional utilizado en la literatura científica y legal para agrupar todas las formas de explotación y abuso sexual infantil facilitadas por tecnologías de la información. \\

    \item[] \textbf{Perfilado Conductual (\textit{Behavioral Profiling}):} Técnica analítica que utiliza algoritmos para identificar patrones en la actividad digital de un usuario a lo largo del tiempo, estableciendo una línea base de comportamiento para detectar desviaciones o anomalías. \\

    \item[] \textbf{Red Neuronal Convolucional (\gls{cnn}):} Arquitectura de aprendizaje profundo (\textit{Deep Learning}) altamente especializada en el procesamiento de datos con topología de cuadrícula, siendo el estándar de la industria para el análisis y clasificación de imágenes en visión artificial. \\

    \item[] \textbf{Servicio de Primer Plano (\textit{Foreground Service}):} Componente crítico del sistema operativo Android que ejecuta operaciones continuas y notorias para el usuario (mediante una notificación fija), evitando que el gestor de memoria cierre la aplicación de monitoreo. \\

    \item[] \textbf{Tasa de Falsos Positivos/Negativos:} Métricas de evaluación de \gls{ia}. Un falso positivo ocurre cuando el sistema alerta de un riesgo inexistente (generando fatiga de alerta), mientras que un falso negativo ocurre cuando el sistema omite una amenaza real. \\

    \item[] \textbf{Vamping:} Comportamiento adolescente caracterizado por el uso prolongado de dispositivos digitales durante la noche y la madrugada, lo que altera severamente los ritmos circadianos y la consolidación de la memoria. \\

    \item[] \textbf{zRAM (Memoria Comprimida):} Tecnología de gestión de memoria a nivel del \textit{kernel} de Linux/Android que comprime bloques de datos inactivos en la memoria \gls{ram} para aumentar el espacio disponible y evitar ralentizaciones durante la ejecución simultánea de aplicaciones. \\
\end{itemize}